% Response to the editor and reviewers. Comments are presented as summaries.
% Build from repository root: make response
% Manuscript page and bibliography references reflect the accompanying revision.
\documentclass[12pt,letterpaper]{article}
\usepackage[margin=1in,headheight=15pt,headsep=18pt]{geometry}
\usepackage{fontspec}
\usepackage{unicode-math}
% Readable submission typography; the journal does not mandate a rebuttal font.
% Keep text and mathematics in the same Times-compatible OpenType family.
\setmainfont{TeX Gyre Termes}
\setmathfont{TeX Gyre Termes Math}
\setmonofont{Latin Modern Mono}[Scale=MatchLowercase]
\usepackage{microtype,amsmath,booktabs,array,longtable,enumitem,needspace}
\usepackage{setspace,ragged2e,fancyhdr}
\usepackage{xcolor}
\usepackage{hyperref}
\usepackage{xurl}
\usepackage{titlesec}
\definecolor{commentink}{HTML}{253F55}
\hypersetup{colorlinks=true,linkcolor=commentink,urlcolor=commentink,
  pdftitle={TRiX response to the editor and reviewers},
  pdfauthor={Stavan Dholakia, Shivani Shukla, Abhishek Singh and Aditya Gazta},
  pdfdisplaydoctitle=true}
\urlstyle{same}
\DeclareUrlCommand\path{\urlstyle{rm}}
\titleformat{\section}{\fontsize{16}{20}\selectfont\bfseries\raggedright}{}{0pt}{}
\titleformat{\subsection}{\fontsize{12}{15}\selectfont\bfseries\raggedright}{}{0pt}{}
\titlespacing*{\section}{0pt}{20pt}{10pt}
\titlespacing*{\subsection}{0pt}{16pt}{6pt}
\setstretch{1.08}
\setlength{\parindent}{0pt}
\setlength{\parskip}{6pt}
\setlength{\emergencystretch}{2em}
\setlength{\RaggedRightParindent}{0pt}
\pagestyle{fancy}
\fancyhf{}
\fancyhead[L]{\small TRiX | Response to reviewers}
\fancyhead[R]{\small Point-by-point response}
\fancyfoot[C]{\small\thepage}
\renewcommand{\headrulewidth}{0.3pt}
\renewcommand{\footrulewidth}{0pt}
\clubpenalty=10000
\widowpenalty=10000
\input{manuscript/data/response_locations.tex}
\newcommand{\mssec}[1]{\csname mssection#1\endcsname}
\newcommand{\msitem}[1]{\csname msitem#1\endcsname}
\newcommand{\mscite}[1]{[\csname mscite#1\endcsname]}
\newcommand{\response}[3]{\Needspace{9\baselineskip}\subsection*{#1\enspace #2}\label{response:#1}{\color{commentink}\itshape\textbf{Comment #1 (summary).} #3\par}\nobreak\textbf{Response.} }
\newcommand{\locations}[1]{\par\nopagebreak[4]{\small\textbf{Manuscript locations:} #1.\par}\nopagebreak[4]}
\newcommand{\evidence}[1]{\par\nopagebreak[4]\begingroup\small\RaggedRight\textbf{Evidence:} \path{#1}\par\endgroup}

\begin{document}
\raggedright
\thispagestyle{plain}
{\fontsize{20}{24}\selectfont\bfseries Response to the editor\\and reviewers\par}
\medskip
{\bfseries\raggedright TRiX: Task-conditioned execution governance\\for robotic disassembly\par}
Manuscript COMMSENG-26-0216-T\par
Stavan Dholakia, Shivani Shukla, Abhishek Singh and Aditya Gazta\par
26 September 2026

\medskip
Comments are summaries. Locations identify the main manuscript or S-prefixed
Supplementary Information, each with separate pagination. Citations use the main
bibliography; \path{archive/} identifies the separate evidence archive.

\section*{Overview of the revision}

For manuscript COMMSENG-26-0216-T, we changed the submitted title,
``TRiX: Neuro-Symbolic Safety for Foundation Model Agents in Robotic
Disassembly'', to ``TRiX: Task-conditioned execution governance for robotic
disassembly''. No autonomous agent loop was evaluated; the model study
measures governance of recorded proposals. The new title makes that scope explicit.

The main changes are:
\begin{enumerate}[leftmargin=*,itemsep=0.3em]
\item Genuine Stable-Baselines3 PPO and SAC experiments with ten training
seeds, recorded selected checkpoints and held-out evaluation. Post-hoc
governance and filter-aware training are reported separately.
\item A matched analytical-command QP control, a learned-representation
decomposition, 200 held-out procedural geometries, and explicit specification
controls. These separate the effects that were conflated in the submission.
\item Four representation and feasibility propositions replace the submitted
ISS theorem. Policy adaptation is reported as an empirical finding.
\item Actual VLM outputs: 273 decisions on 91 shared PCB images. Of 257
grounded command pairs, TRiX prevents all 31 recorded constraint-violation
cases and introduces none during the 40-step horizon.
\item Physical B601-RS intervention evidence and three video-documented
grasp-and-return trials, plus three additional logged repetitions.
\item Mapped source and evidence, seed-level uncertainty, disclosed provenance
limits, and code-generated workspaces driven by recorded environment states.
\end{enumerate}

\clearpage
\section*{Response to the Editor}

\response{E.1}{Positioning among existing safety methods}{Expand the discussion of standard safety mechanisms beyond Lagrangian methods.}
The introduction and comparison table distinguish constrained
policy optimization, projection-based intervention, recovery policies and
control barrier functions. Table~1 now describes the cited formulations and
their assumptions, rather than assigning categorical capabilities to entire
method classes. The comparison separates learning a task policy, representing
a constraint and enforcing a command at runtime. We explicitly acknowledge
that CBF methods can encode nonlinear and state-dependent constraints.
TRiX's contribution is stated in terms of task-conditioned execution authority
and its interaction with the upstream proposal distribution.
\locations{\mssec{1}; \mssec{2.1}; \mssec{4.3}; \msitem{tab:submitted-1}; manuscript references \mscite{ref16}--\mscite{ref26}}
\evidence{manuscript/TRIX_REVISION.tex}

\response{E.2}{Introduce the architecture figure in context}{Reference Figure~1 in the text and place it where its purpose is clear.}
Figure~1 now appears with the foundation-model motivation and is explicitly
introduced in the prose. It shows the proposal, task state/context, TRiX's
decision and the execution boundary. ALLOW and PROJECT can reach execution;
REJECT and INFEASIBLE terminate without an execution command. This makes the
architecture available before the formal definitions.
\locations{\mssec{1}; \mssec{2.1}; \msitem{fig:submitted-1}}
\evidence{manuscript/figures/trix_architecture_fig1.svg}

\response{E.3}{Safety, damage and task quality}{Establish which failures are safety-relevant and distinguish them from quality or task-performance failures.}
We distinguish modeled command inadmissibility, outcome-affecting violations,
diagnostic violations, mechanical failure and successful completion. A violated
command boundary is not automatically a damaged workpiece. BATTERY provides
a hazard-motivated example through its declared deformation, internal-short
and delayed-heating model; its physical thresholds have not been independently
validated on hardware. PCB fracture, latch breakage and thread stripping are
asset-integrity outcomes. Progress or timeout measures task performance.
The episode contract and constraint table specify which recorded events
change the evaluation outcome. The VLM result is correspondingly reported
as prevention of constraint-violation cases, without claiming prevented fractures.
\locations{\mssec{3}; \mssec{4.1}; \mssec{S1}; \mssec{S4.4}; \mssec{S13}; \msitem{tab:submitted-8}}
\evidence{benchmark/outcome.py; envs/battery_env_v2.py; manuscript/TRIX_SUPPLEMENT.tex}

\response{E.4}{Quantitative latency}{Replace vague claims of negligible latency with a quantitative computational comparison.}
We replaced the unsupported historical timing figures with a new recorded
benchmark of 22 existing filter implementations across all six tasks.
Five fresh processes each evaluate the same 2,048 saved inputs per task,
following 256 warm-up calls per implementation. The resulting 225,280 timed
calls, output actions, failure flags, execution order, host information and
source hashes are supplied with the reproducible benchmark.

For the matched SCREW command polytope, median filter-call latency is
1.255~$\mu$s for the production scalar TRiX projector,
12.578~$\mu$s for active-set QP and 14.924~$\mu$s for reused OSQP.
The corresponding 95th percentiles are 1.556, 15.534 and 19.895~$\mu$s.
All matched implementations pass the declared output-agreement and
constraint-residual checks with zero recovery outputs. The maximum
action-coordinate difference is $1.013\times10^{-6}$ on the saved bank.
\msitem{tab:runtime-latency} reports every implementation, including
the NumPy analytical projector, whose 18.163~$\mu$s median shows that
implementation choices materially affect cost.

The timer covers the complete filter callable, including allocation,
per-call constraint construction and diagnostics. Policy/VLM inference,
simulation and robot communication are excluded. Supplementary Section~S8 specifies
the Intel Core Ultra 9 275HX host, numerical package versions, CPU affinity,
solver tolerances, setup reuse and the full sampling protocol. These are
measured implementation costs on the recorded host, not a claim of
worst-case execution time or end-to-end robot latency.
\locations{\mssec{2.7}; \mssec{4.8}; \mssec{S8}; \msitem{tab:runtime-latency}}
\evidence{experiments/runtime_benchmark.py; docs/runtime_benchmark_protocol.md; results/runtime/manifest.json; results/runtime/summary.json; manuscript/TRIX_SUPPLEMENT.tex}

\Needspace{23\baselineskip}
\response{E.5}{Generality beyond six task implementations}{Explain how the method generalizes beyond the six manually specified disassembly tasks.}
We define TRiX through a task-conditioned admissible set and four execution
outcomes. The task-specific work is to identify the action units, constraint
model, required state/context, uncertainty allowance and admissible correction
or rejection rule. The six tasks instantiate this construction. A separate
procedural benchmark evaluates box, ellipsoidal, polytope, state-scaled,
rotating and phase-gated sets on 200 held-out instances. This extends the
representation analysis beyond named task examples.

The protocol is explicit: learned constraint models are fitted separately
on each evaluated geometry, and the context-free box defines the structural
partition. This is neither zero-shot transfer of one learned model nor an
independent discovery of that partition. The preregistered BAYONET test is
also retained as inconclusive because no learned arm met its competence gate.
Together, these results identify where the formulation has been tested and
where further task modeling and evaluation remain necessary.
\locations{\mssec{2.1}; \mssec{2.4}; \mssec{3}; \mssec{4.4}; \mssec{S10}; \msitem{tab:procedural-geometry}}
\evidence{benchmark/geometry_gen.py; results/geometry_bench_heldout.json; docs/bayonet_preregistration.md; manuscript/TRIX_SUPPLEMENT.tex}

\clearpage
\section*{Response to Reviewer 1}

\response{R1.1}{References in comparison tables}{Add representative references for the paradigms in Tables~1 and~3.}
Both tables now identify the literature or implementation supporting each
entry. Table~1 cites constrained-policy, safety-layer, recovery and CBF
formulations. Table~3 separates trained PPO/SAC policies from analytical
runtime restrictions, reference controls and the learned-constraint study.
The latter is not represented as a trained SafeLayer task policy. This avoids
implying an experimental baseline merely because its method family appears
in related work.
\locations{\mssec{4.3}; \msitem{tab:submitted-1}; \msitem{tab:submitted-3}}
\evidence{benchmark/registry.py; training/sb3_runner.py}

\response{R1.2}{Accessibility of the theoretical framework}{Add background and make the methodology easier for non-specialist readers to follow.}
The exposition now starts with a concrete execution question: whether a
proposed command can be forwarded, corrected or rejected. The architecture
figure precedes the mathematics. The screw example introduces coupled
translation and rotation in normalized action coordinates; the PCB example
then introduces a coupled inequality, followed by phase and coordinate-frame
dependence. The main text summarizes four propositions, with the proofs and
practical interpretations in Supplementary Sections~S1--S3. The theory's scope is stated before the experiments, and the
former broad ISS argument is removed.
\locations{\mssec{1}; \mssec{2.1}; \mssec{S1}; \mssec{S2}; \mssec{S3}; \msitem{fig:submitted-1}; \msitem{fig:submitted-2}}
\evidence{manuscript/TRIX_REVISION.tex; manuscript/TRIX_SUPPLEMENT.tex}

\response{R1.3}{Simulated workspace and interaction}{Show the simulated workspaces with an end-effector interacting with task-specific workpieces.}
Figure~3 now shows all six task workspaces with the B601 gripper, task
workpiece, coordinate grid and applied-command cues. The panels are
deterministic PyBullet visualizations driven by recorded analytical
environment states. The state manifest records each selected frame, state
value and source mapping, and the six renderer entrypoints and dependencies
are included in the source reference.

The caption and Methods distinguish the two layers of the implementation:
the analytical environments generate the dynamics and reported outcomes;
PyBullet visualizes their archived states. The figure does not imply that
articulated-arm or contact dynamics were used to generate the policy results.
\locations{\mssec{2.2}; \mssec{4.1}; \mssec{S4}; \msitem{fig:disasm-workspaces}}
\evidence{manuscript/figures/workspaces/state_manifest.json; docs/workspace_rendering.md; assets/workspaces/frozen_frames.json; manuscript/TRIX_SUPPLEMENT.tex}

\response{R1.4}{Concrete model proposals and corrections}{Provide concrete examples of foundation-model errors and how TRiX corrects them.}
We replaced the synthetic injection experiment with actual outputs from three
provider runs on a shared image bank. A specific recorded GPT-4o response to
\texttt{s06\_occluded.png} proposes \texttt{LIFT} at magnitude 0.8, which
grounds to normalized command $(0,0,0.8)$. The archived unfiltered replay has
40 violating steps out of 40; its governed counterpart has zero. Across the
257 grounded command pairs, governance prevents all 31 observed violation
cases and introduces none. A changed command alone receives no prevention credit.

The manuscript distinguishes inadmissibility from hallucination, refusal and
ungroundable output. These records contain no independent human annotation
of semantic correctness. The example demonstrates command governance over
the fixed replay horizon, not successful autonomous extraction. R4.9 below
gives the complete numerical and visual-provenance account.
\locations{\mssec{2.5}; \mssec{4.5}; \mssec{S11}; \msitem{tab:vlm-results}}
\evidence{results/vlm/vlm_openai.json; manuscript/data/vlm_trial_ledger.json; manuscript/TRIX_SUPPLEMENT.tex}

\response{R1.5}{Adding or changing constraints}{Explain scalability to additional tasks or modified physical constraints.}
The revised formulation makes the task-dependent inputs explicit:
$\mathcal A(x,c)$, command scaling, state/context selection, uncertainty
tightening and the solution or rejection rule. A changed constraint must be
represented and its feasibility checked; projection cannot supply a missing
physical specification. The procedural experiment tests six set families on
200 held-out instances, while the sensitivity experiment varies 21 constants
over five multipliers. These address representation breadth and specification
dependence, respectively. They do not establish computational scaling with
arbitrary action dimension, automatic constraint discovery or transfer to an
unmodeled manipulation task.
\locations{\mssec{2.1}; \mssec{2.4}; \mssec{3}; \mssec{4.4}; \mssec{S10}}
\evidence{benchmark/geometry_gen.py; results/geometry_bench_heldout.json; results/sensitivity.json; manuscript/TRIX_SUPPLEMENT.tex}

\response{R1.6}{Computational time}{Report the computation time required by TRiX and the comparator.}
We added the recorded computational comparison in Supplementary Section~S8. Across the six
registered TRiX task arms, median complete-filter latency ranges from
1.123~$\mu$s for PCB to 11.770~$\mu$s for CRANK; BATTERY's preventive
variant takes 3.122~$\mu$s. The matched SCREW comparison measures
1.255~$\mu$s for production TRiX, 12.578~$\mu$s for active-set QP and
14.924~$\mu$s for reused OSQP. The table reports medians, 95th percentiles
and variation across five fresh processes, with all raw calls retained.

Each row uses 10,240 calls on the same saved input bank for its task.
The source specifies solver settings, warm-up, a single pinned CPU and the
measurement boundary. OSQP's one-time object construction is reported
separately from steady-state calls. The different componentwise and
oracle-tangent restrictions are identified explicitly, so their timings
are not presented as equivalent-safety comparisons. Response E.4 gives
the matched-solver agreement checks and the timing scope.
\locations{\mssec{2.7}; \mssec{4.8}; \mssec{S8}; \msitem{tab:runtime-latency}; response E.4}
\evidence{experiments/runtime_benchmark.py; docs/runtime_benchmark_protocol.md; results/runtime/manifest.json; results/runtime/summary.json; manuscript/TRIX_SUPPLEMENT.tex}

\response{R1.7}{Distributions and uncertainty}{Show distributions or boxplots rather than relying only on aggregate performance.}
Figure~4 shows all ten training-seed values for each displayed arm and
regime, with the underlying 80 values retained in the source package.
The plotted results expose the CRANK seed heterogeneity and the small PRY
mean difference, alongside the PCB regime change and SNAP comparison.
\msitem{tab:submitted-11} adds 95\% percentile bootstrap intervals for 13 selected seed-paired
contrasts. Training seed is the experimental unit; the 100 evaluation episodes
within each seed are not treated as independent trained-policy replicates.
The table identifies comparisons that use separately trained policies or
different episode streams. A collapsed interval is not an equivalence test.
\locations{\mssec{2.3}; \mssec{4.6}; \mssec{S6}; \msitem{fig:seed-distributions}; \msitem{tab:submitted-11}}
\evidence{manuscript/figures/figure4_seed_data.csv; manuscript/data/seed_statistics_pairs.csv; scripts/build_revision_seed_statistics.py; manuscript/TRIX_SUPPLEMENT.tex}

\response{R1.8}{Figure references and Figure~4 quality}{Improve figure integration and the presentation of Figure~4.}
All five figures now have explicit callouts in the body. Figure~4 replaces
the submitted learning-curve plot from a different environment with seed-level
results from the revised production records. Its caption identifies the
post-hoc and filter-aware regimes, trained algorithm and episode count. The
saved 80 plotted values reconcile with the corresponding revised results.
Figure~4 has been regenerated from the retained 80 seed-level values using
\path{scripts/build_revision_figure4.py}. This replacement builder checks every
value against the retained audited seed counts before plotting. The historical
plotting script remains unavailable; the replacement and its input/output hashes
are supplied with the source.
\locations{\mssec{1}; \mssec{2.1}; \mssec{2.2}; \mssec{2.3}; \mssec{2.6}; \msitem{fig:seed-distributions}}
\evidence{manuscript/figures/figure4_seed_distributions.pdf; manuscript/figures/figure4_seed_data.csv; scripts/build_revision_figure4.py; manuscript/figures/figure4_rebuild.json; docs/REPRODUCING.md}

\response{R1.9}{Foundation-model robotics literature}{Expand discussion of LLM/VLM methods in collaborative and embodied robotics.}
The introduction cites PaLM-E, RT-2, SayCan, Code as Policies and
ChatGPT for Robotics, together with Rrapi et al.'s 2026 survey of AI and
language models in collaborative robotics (DOI: 10.1016/j.rcim.2026.103269;
PII: S0736584526000487), to motivate the expanding range of learned and
generative action proposals. The survey is cited in context for its discussion
of reasoning, planning and interaction. This literature motivates the
separate execution-authority question, rather than treating semantic planning
as physical admissibility. The actual VLM study then evaluates that interface
through recorded proposals and paired command execution. The paper does not
claim a human-robot collaboration experiment.
\locations{\mssec{1}; \mssec{2.5}; manuscript references \mscite{ref8}--\mscite{ref13}, \mscite{ref39}}
\evidence{manuscript/TRIX_REVISION.tex}

\clearpage
\section*{Response to Reviewer 2}

\response{R2.1}{PCB variables and physical interpretation}{Define the PCB variables and explain the physical meaning of the constraint.}
The revised PCB interface specifies two commanded tilt torques,
$\tau_x$ and $\tau_y$, and an axial lift force $F_z$. These are distinguished
from the state variables: two tilt angles, angular rates, lift and lift velocity.
The nominal torque condition is
$\sqrt{\tau_x^2+\tau_y^2}\leq0.100$ N\,m, not a zero-torque requirement.
The implemented robust projector uses a smaller radius, and that difference
is disclosed. Table~2 summarizes the modeled commands and damage conditions;
Supplementary Section~S4 gives the complete command scales, observation
indices and units. Supplementary Section~S2.2 supplies
the disk--interval projection.
\locations{\mssec{2.2}; \mssec{4.1}; \mssec{S2.2}; \mssec{S4.1}; \mssec{S4.2}; \msitem{tab:submitted-2}; \msitem{tab:action-scales}}
\evidence{envs/pcb_env_v2.py; baselines/pcb_filters.py; manuscript/TRIX_SUPPLEMENT.tex}

\response{R2.2}{Observation dimensionality}{Resolve the inconsistent state/observation dimension and provide its components.}
All six revised Gymnasium observations have ten entries. Their meanings are
task-specific; they are not a shared position/velocity/wrench vector. \msitem{tab:submitted-7}
lists every populated index, physical meaning, derived flag and zero-padded
entry for each environment. For example, PCB contains lift and lift velocity,
two tilt angles, two angular rates, accumulated damage and a fracture flag,
followed by two zeros. Reset-sampled parameters are not extra observation
coordinates. Context-limited filters receive the same observation but use
fewer of its fields. The submitted dimensional description is superseded
by this code-derived table.
\locations{\mssec{4.1}; \mssec{S4.1}; \msitem{tab:submitted-7}}
\evidence{benchmark/gym_adapter.py; envs/pcb_env_v2.py; envs/snap_env_v2.py; envs/crank_env_v2.py; manuscript/TRIX_SUPPLEMENT.tex}

\response{R2.3}{Reward symbols and outcome notation}{Define the progress and violation terms and use consistent notation.}
We removed the submitted universal reward expression and its undefined
$v_{\mathrm{progress}}$ and $I_{\mathrm{violation}}$ symbols. The revised
description follows the task-specific implemented rewards and identifies
the normalized progress quantity used for checkpoint selection: extraction,
lift, rotation, peel or gap opening as applicable. Reward for learning is
separated from the evaluation outcome contract. The latter partitions episodes
into five mutually exclusive outcomes and distinguishes outcome-affecting
from diagnostic violations. Checkpoint selection ranks safe completion
before progress, preventing a stationary policy from winning solely by
avoiding a violation.
\locations{\mssec{4.1}; \mssec{4.2}; \mssec{S4.3}; \mssec{S4.4}}
\evidence{benchmark/outcome.py; training/stage1.py; envs/pcb_env_v2.py; manuscript/TRIX_SUPPLEMENT.tex}

\response{R2.4}{Domain randomisation appendix}{Fill the empty appendix with parameters, distributions, ranges and training versus out-of-distribution use.}
The requested systematic table is now \msitem{tab:submitted-9} in
Supplementary Section~S4.6. It lists each reset parameter, distribution,
numerical range and training/validation/test use, separately from per-step
Gaussian noise. The original held-out evaluation uses the training ranges.

To address the OOD-use request with direct policy evidence, we added a new
experiment during revision. Its declaration and implementation were committed
as \texttt{9bc2242} before any OOD rollout. We evaluated the seven eligible
Table~4 comparisons across five tasks, both arms, ten training seeds and
100 episodes per cell under three declared conditions: $1.5\times$ and
$2\times$ reset half-widths and a $2\times$ outside-support shell. All 420
cells and 42,000 episodes are reported. The 110 selected checkpoints were
neither retrained nor reselected. Governor definitions, margins, thresholds,
noise and episode caps remained fixed, with matching constraint, margin and
taxonomy hashes and separately recorded evaluation-distribution hashes.
SCREW uses a jointly sampled friction pair conditioned on $\mu_s\geq\mu_k$.
PRY is marked not applicable because its physical reset is deterministic;
we did not introduce a new bond-strength randomisation.

\msitem{tab:ood-ranges} gives the added ranges; \msitem{tab:ood-safe} reports
absolute safe completion, changes from historical in-range values and exact
seed-bootstrap intervals for every condition. \msitem{tab:ood-outcomes}
reports destructive, unsafe intact, mechanical-failure, timeout and violation
counts, plus within/outside-support denominators. The same training-seed
estimator is used as in main Section~4.6; historical comparisons are not
episode-paired and therefore include episode-sampling variation.

In SNAP's outside-support shell, TRiX safe completion falls from 96.5\% to
59.9\%, versus 3.5\% under static, while fractures occur in 401 versus 468
of 1,000 episodes per arm, including timeouts. Destructive completions
number 369 versus 886; this separate category pools completed episodes with
fracture or accumulated damage and excludes timeouts. The larger advantage
in safe completion therefore accompanies a smaller reduction in fractures.
Both measures are stated together in the main Results and Discussion, with
the full partition in \msitem{tab:snap-ood-attribution}.

A post-hoc audit exactly reproduced all 6,000 SNAP episodes and found all
293,082 TRiX commands inside the frozen executable bounds. All TRiX shell
fractures began through overdeflection in the softer-latch stratum, before
pull damage. \mssec{S14.4} also discloses that the historical executed
margins differ from the later, tighter registry, and quantifies that
discrepancy without changing the projector. This result distinguishes
command membership from episode safety. The use column now records both
native-range evaluation and the added OOD conditions; original in-range
results remain unchanged.

The separate 105-case sensitivity sweep in Supplementary Section~S7 remains
a scripted study of physical/specification constants, with 96/105 cases
meeting its prescribed ordering. It is not pooled with this new reset-distribution
evaluation and does not substitute for it. Neither study establishes independent
physical calibration or generalisation to arbitrary unseen dynamics.
\locations{\mssec{2.3}; \mssec{3}; \mssec{4.1}; \mssec{4.6}; \mssec{S4.6}; \mssec{S14}; \mssec{S14.4}; \msitem{tab:submitted-9}; \msitem{tab:ood-ranges}; \msitem{tab:ood-safe}; \msitem{tab:ood-outcomes}; \msitem{tab:snap-ood-attribution}}
\evidence{docs/ood_reset_protocol.md; docs/ood_reset_protocol.json; experiments/ood_reset.py; benchmark/evaluation_distribution.py; training/sb3_runner.py; results/ood_reset/summary.json; results/ood_reset/seed_records.json; results/ood_reset/reset_draws.json; results/ood_reset/manifest.json; docs/ood_reset_results.md; docs/snap_ood_attribution_plan.md; docs/snap_ood_attribution_results.md; experiments/snap_ood_attribution.py; results/snap_ood_attribution/summary.json; results/snap_ood_attribution/episode_attribution.json; results/snap_ood_attribution/verification.json; results/snap_ood_attribution/manifest.json; experiments/sensitivity.py; results/sensitivity.json}

\response{R2.5}{Placement of runtime discussion}{Move computational timing out of the theoretical derivation.}
The helical-coupling subsection now contains only the model, normalization
and projection interpretation. Computational scope is treated separately in
Supplementary Section~S8 and the solver-control Methods subsection. The new timing protocol
and measured distributions are reported there, with their raw evidence and
implementation scope. They replace the unsupported historical figures.
Responses E.4 and R1.6 describe the completed measurement.
\locations{\mssec{2.1}; \mssec{2.7}; \mssec{4.8}; \mssec{S8}; \msitem{tab:runtime-latency}}
\evidence{manuscript/TRIX_REVISION.tex; results/runtime/summary.json; manuscript/TRIX_SUPPLEMENT.tex}

\response{R2.6}{Orphaned Figure~2 reference}{Repair the disconnected sentence referring to Figure~2.}
Figure~2 is integrated with the screw-coupling explanation and explicitly
referenced from the canonical example. Its caption defines the normalized
axes, physical command scales, finite tolerance and the status of the
illustrative componentwise box. The figure therefore supports the local
argument about coupled commands, with no detached reference to a different
experiment.
\locations{\mssec{2.1}; \msitem{fig:submitted-2}}
\evidence{manuscript/figures/figure2_geometry.json; scripts/build_revision_figure2.py}

\response{R2.7}{Narrative structure}{Connect the physical constraint, the limitation of the compared representation and the resulting TRiX design.}
We reorganized the argument around that sequence. Screw coupling first
shows what independent coordinate bounds omit. PCB then illustrates a coupled
norm, while phase and frame examples explain which state information selects
the applicable set. The formal propositions identify those representation
properties; the experiments test their consequences under post-hoc governance
and filter-aware training. The discussion includes controls where simple
restrictions suffice and cases where they remain limiting. This links the
physical requirement to the implemented comparator without presenting a
universal ranking of safety algorithms.
\locations{\mssec{1}; \mssec{2.1}; \mssec{2.3}; \mssec{2.4}; \mssec{3}}
\evidence{manuscript/TRIX_REVISION.tex}

\clearpage
\section*{Response to Reviewer 4}

\textbf{Title change.} For manuscript COMMSENG-26-0216-T, we changed the submitted title,
``TRiX: Neuro-Symbolic Safety for Foundation Model Agents in Robotic
Disassembly'', to ``TRiX: Task-conditioned execution governance for robotic
disassembly''. No autonomous agent loop was evaluated; the model study
measures governance of recorded proposals. The new title makes that scope explicit.


\response{R4.1}{Trained policies and experimental provenance}{The submitted results used random-action transformations rather than the claimed trained policies.}
The criticism identifies a substantive defect in the submitted evidence.
We replaced those headline results with recorded Stable-Baselines3 PPO and
SAC experiments. The revision claims only those two trained policy algorithms;
PPO-Lagrangian and a trained SafeLayer task policy are not present in the
verified production matrix. Learned constraint representations are reported
separately.

Methods specifies ten production training seeds, checkpoint candidates,
validation ranking and held-out evaluation. Stage~1 applies runtime arms to
frozen nominal policies; Stage~2 trains with the restriction active. Main
Table~4 summarizes the central comparisons; \msitem{tab:submitted-4}
reports the full Stage~1 matrix and \msitem{tab:stage2-key} the Stage~2 summary. The release contains 520 selected
policy files, the candidate checkpoints and the evaluation shards. The full
audit links 1,020 recorded aggregates to 20,400 exact-arm test shards.

The audit also found a BATTERY filename-prefix collision that pooled distinct
arms in 30 historical aggregates. Counts were recomputed from the shards of
the recorded selected checkpoints. The Stage~1 SAC/TRiX value in \msitem{tab:submitted-4} is 20.0\%,
replacing the intermediate draft's 59.8\%. The exact-arm audit is
\path{docs/battery_aggregation_correction.md}; the archive's
\path{verification/policy_record_linkage.json} records the corresponding
\texttt{stage1/BATTERY/sac/trix} group and each contributing shard.
The original files and source matcher correction are retained. No
checkpoint was reselected or rerun; this correction does not establish that
repeating historical selection would choose the same checkpoint. The separate
PCB oracle seed correction is also retained. Seed-level statistical inference
replaces the submitted episode-stream tests.
\locations{\mssec{4.2}; \mssec{4.6}; \mssec{S5}; \mssec{S6}; \mssec{S9}; \msitem{tab:submitted-4}; \msitem{tab:stage2-key}}
\evidence{training/sb3_runner.py; docs/battery_aggregation_correction.md; archive/verification/policy_record_linkage.json; manuscript/TRIX_SUPPLEMENT.tex}

\response{R4.2}{Matched analytical comparator}{The SafeLayer comparison is unfair; compare against a CBF-QP supplied with the same analytical constraint.}
We removed the inference that a learned approximation's failure demonstrates
inferiority of classical constraint enforcement. The revised SCREW control
supplies the analytical projector and registered \texttt{qp\_matched} arm
with the same finite-tolerance command polytope, bounds, margins and
normalized Euclidean objective. Under post-hoc evaluation, the matched
analytical reference, QP and TRiX each achieve 88.7\% safe completion for
SAC and 60.0\% for PPO. With the restriction active during training, the
matched constrained arms reach 100\% for both algorithms.

The implemented production comparator is an analytical command-set QP.
We label it \emph{matched QP}, because equality of that optimization problem
does not constitute a separate state-space barrier synthesis or proof of
forward invariance. This directly tests whether TRiX's projection machinery
confers an advantage when the specification is matched: no task-level
advantage is observed here. We do not claim a full dynamics-aware CBF study
across all six tasks.
\locations{\mssec{2.1}; \mssec{2.3}; \mssec{4.3}; \mssec{S2.1}; \mssec{S8}; \msitem{tab:submitted-4}}
\evidence{baselines/screw_solvers.py; baselines/screw_filters.py; results/stage1/stage1_records.json; results/stage2/stage2_records.json; manuscript/TRIX_SUPPLEMENT.tex}

\response{R4.3}{CBF capabilities}{Correct Table~1's categorical statements about control barrier functions.}
The categorical exclusions have been removed. The text now recognizes
nonlinear safe sets, state-dependent barriers and extensions involving
relative degree and time variation. Table~1 describes representative
formulations with citations and explicit assumptions. The manuscript separates
command-set membership from state-space forward invariance and states that
barrier synthesis requires its own dynamical conditions. TRiX is positioned
through its execution contract and empirical regime comparisons.
\locations{\mssec{1}; \mssec{4.3}; \msitem{tab:submitted-1}; manuscript references \mscite{ref23}--\mscite{ref25}}
\evidence{manuscript/TRIX_REVISION.tex}

\response{R4.4}{ISS theorem assumptions}{The submitted stability assumptions substantially supply the claimed conclusion.}
We withdraw the submitted ISS theorem and its associated guarantee. The
revision instead proves four specific statements: projection onto a product
of intervals equals clipping; a phase-independent restriction safe in two
phases must lie in their intersection; an orientation-independent restriction
can lose authority for rotating anisotropic sets; and Gaussian tightening can
remove all nonzero or all feasible command authority. Each statement gives
its assumptions, proof and interpretation. These are representation and
feasibility results, not closed-loop stability or physical-safety theorems.
Policy adaptation is explicitly empirical and is not a fifth proposition.
\locations{\mssec{2.1}; \mssec{3}; \mssec{S3}}
\evidence{manuscript/TRIX_REVISION.tex; manuscript/TRIX_SUPPLEMENT.tex}

\response{R4.5}{Non-expansiveness of projection}{The projection argument does not hold for arbitrary differentiable manifolds.}
We agree and remove the unrestricted non-expansiveness claim and the proof
that relied on it. The revised global projection definition applies to
nonempty closed convex sets with normalized command coordinates; the affine
case has the stated full-row-rank condition. General admissible sets may
be nonconvex, and the manuscript no longer derives global uniqueness,
non-expansiveness or stability from differentiability alone. Phase changes
and frame selection are treated as context-dependent set selection, with
the limitations of the resulting execution guarantee stated separately.
\locations{\mssec{2.1}; \mssec{S1}; \mssec{S2}; \mssec{S3}}
\evidence{manuscript/TRIX_REVISION.tex; manuscript/TRIX_SUPPLEMENT.tex}

\response{R4.6}{Circular validation}{Agreement between the NumPy environment and constraints derived from its own equations is not independent physical validation.}
We agree. The revised Methods explicitly calls the benchmark analytical and
withdraws the submitted same-model/PyBullet validation argument. The screw
relation is presented as an ideal kinematic identity. Acceptance tests verify
that the implementation follows the declared equations; they do not verify
the physical fidelity of those equations or their numerical parameters.

We separate this consistency check from two further analyses. First, matched
SCREW operators share a declared optimization problem, and the archived
policy results establish matching task-level outcomes. The new finite-bank
agreement checks in Supplementary Section~S8 compare analytical, active-set and
operator-splitting implementations; agreement does not validate the physical
model. Second, the sensitivity sweep tests
21 constants at five multipliers with eight scripted episodes per arm and
case. The prescribed ordering holds in 96 of 105 cases. The nine exceptions
comprise six PCB model/boundary inconsistencies, one CRANK tightening failure
and two BATTERY cases with zero destructive completions but lower safe
completion for the designated preventive arm. This is a stratified empirical
diagnosis, not a 96/96 safety guarantee after excluding failures.

The task definitions, source fingerprints and assumptions make the model
auditable, while the physical calibration of the benchmark force, torque and
damage thresholds remains unestablished. Neither the VLM study nor the
separate B601 execution demonstration is offered as independent validation
of those quantities.
\locations{\mssec{2.1}; \mssec{3}; \mssec{4.4}; \mssec{S4.4}; \mssec{S7}; \mssec{S8}}
\evidence{results/sensitivity.json; experiments/sensitivity.py; baselines/screw_solvers.py; assets/workspaces/environment_sources.json; results/runtime/summary.json; manuscript/TRIX_SUPPLEMENT.tex}

\response{R4.7}{Role of PyBullet}{Explain the PyBullet code in the repository and its relationship to the reported experiments.}
The production experiments use the versioned analytical environments.
An earlier helical-constraint prototype, retained as
\path{research/experiments/pybullet_helix.py} in the evidence archive,
imposed the analytical screw pitch through a gear constraint. It therefore
could not independently validate that pitch. We withdrew the associated
validation argument; no result from that prototype is used in this revision.
PyBullet supplies the deterministic visualization layer for the six workspace
panels. The manuscript states this separation and cites the rendering
implementation. The supplied per-panel manifests bind workpiece poses to
archived environment states. We do not present these renders as a second
independent physics validation, and the illustrated gripper does not imply
that articulated-arm or contact dynamics generated the reported task outcomes.
\locations{\mssec{2.2}; \mssec{4.1}; \mssec{S4}; \msitem{fig:disasm-workspaces}; manuscript reference \mscite{ref33}}
\evidence{docs/workspace_rendering.md; manuscript/figures/workspaces/state_manifest.json; scripts/workspace_environment.py; manuscript/TRIX_SUPPLEMENT.tex; archive/research/experiments/pybullet_helix.py}

\response{R4.8}{Physical robot evidence}{The submission contains no hardware demonstration.}
We added physical execution evidence on the Seeed reBot B601-RS. The archived
intervention experiment records an invalid symbolic transition and a request
to move protected Joint~1 being rejected before physical target generation.
An excessive Joint~2 proposal of $+50^\circ$ is projected to $+18^\circ$;
only the projected target is sent to and executed by the robot. The accepted
return finishes with a recorded Joint~1 drift of $+0.007^\circ$.

The frozen final implementation then completes three continuous-video
grasp-and-return trials with the object retained. Each has ten ALLOW decisions;
these trials are not counted as additional PROJECT/REJECT interventions.
The largest absolute Joint~1 return drift is $0.027^\circ$, and the largest
final residual across Joints~1--6 is $0.122^\circ$. Three further logged
repetitions completed normally; they are not counted as video-verified
retention trials. The archive preserves logs, passive CAN records, source
snapshots, checksums and video mappings for both evidence groups.

This establishes physical deployment of the execution boundary and the
recorded intervention behavior. There is no independent wrench sensor, and
the hardware study does not validate the six benchmark damage models or
provide disassembly-level physical safety statistics.
\locations{\mssec{2.6}; \mssec{3}; \mssec{4.7}; \mssec{S12}; \msitem{tab:b601-hardware}; \msitem{fig:b601-hardware}}
\evidence{archive/hardware/experiments/hardware/2026-08-27_trix_physical_safety_v1; archive/hardware/results/hardware/v4_confirmation/final_submission/HARDWARE_EVIDENCE_README.md; manuscript/TRIX_SUPPLEMENT.tex}

\response{R4.9}{Actual VLM evaluation}{Synthetic hallucination injection does not establish governance of actual VLM outputs.}
We replaced the synthetic injection experiment with actual model responses.
Runs configured as Claude Sonnet, GPT-4o and Gemini each received the same
91 PCB images and prompt, producing 273 saved decisions. A fixed parser
grounds proposals into commands; each grounded command is evaluated with
and without TRiX from the same restored state and seed. The primary result
is prevention of all 31 observed constraint-violation cases, with zero new
cases among the 257 grounded pairs during the 40-step horizon.
Results Section~2.5 now states alongside that count that the commands were
elicited from visibility-scaled images whose displayed tilts can exceed the
intact-board range, and that the result concerns the recorded commands.

The decision accounting is complete:
\begin{center}\small
\begin{tabular}{lrrrrr}
\toprule
Run & Refused & Admissible & Inadmissible & Grounded & Prevented\\
\midrule
Claude Sonnet & 0 & 88 & 3 & 91 & 0\\
GPT-4o & 16 & 66 & 9 & 75 & 9\\
Gemini & 0 & 65 & 26 & 91 & 22\\
\midrule
Total & 16 & 219 & 38 & 257 & 31\\
\bottomrule
\end{tabular}\end{center}
No ungroundable output remains. Gemini's admissible count is 65, correcting
the intermediate draft's 62. Both admissible and inadmissible grounded
commands are replayable. The three inadmissible Claude proposals have no
unfiltered violation and receive no prevention credit. Refusals remain in
the all-image denominator but are not replayed.

For a concrete example, GPT-4o's saved \texttt{s06\_occluded.png} response
grounds to $(0,0,0.8)$. The unfiltered record has 40 violating steps out of
40, versus zero out of 40 under TRiX. Across the bank, 35 images elicit an
inadmissible proposal in at least one run and none across all three. The
Claude and GPT-4o inadmissible sets are disjoint by both image and base-state
ID; holding those responses fixed, the no-three-way-overlap statement holds
for any Gemini failure set.

\textbf{Visual encoding.} The bank contains 28 base configurations:
21 with four visual conditions each and seven partial views. Main-stratum
display tilt is approximately 20.669 times the mechanical replay curvature;
partial views use 3.6 times. The display is a visibility-scaled cue, not a
physically realizable intact-board pose under the benchmark model. Saved
near-boundary cases display $42.4^\circ$--$50.4^\circ$, whereas the model
fractures above $0.075$ rad ($4.30^\circ$) of mechanical curvature. The
positive scale preserves main-stratum ordering, not physical angular distance
to a boundary. All margin and visual-tilt comparisons use only the 84
main-stratum images; the seven partial views remain separate. Rendering
reproducibility does not establish identity of displayed and replayed tilt.
The stratum counts describe responses to this encoding. Without an unscaled
image control, the contribution of visual scaling cannot be separated from
differences between model runs.

\textbf{Model provenance.} The recorded strings are
\texttt{claude-sonnet-4-5}, \texttt{gpt-4o} and
\texttt{gemini-flash-latest}. They are configured aliases, not pinned backend
snapshots. The archived client allows fallback; resolved backend IDs,
API access times and attempt-level histories were not retained. Earlier
notes report four retried Gemini timeouts, but their image IDs and successful
variants cannot be reconstructed. The ledger marks these fields unknown.
If at most four Gemini classifications changed, its inadmissible count would
be 22--30, still above the other runs' fixed counts of nine and three. This
is a conditional sensitivity calculation: the retry note does not establish
that only four decisions could have used another backend.

The paired outcome concerns violations of the declared command constraint.
No replay records mechanical failure or successful extraction; this is a
single-proposal governance study, not an autonomous agent evaluation. The
result remains prominent in the abstract, Results and conclusion. Its
measurement and provenance boundaries are stated alongside it.
\locations{\mssec{2.5}; \mssec{4.5}; \mssec{4.6}; \mssec{S11}; \msitem{tab:vlm-results}; manuscript references \mscite{ref36}--\mscite{ref38}}
\evidence{results/vlm/vlm_anthropic.json; results/vlm/vlm_openai.json; results/vlm/vlm_google.json; manuscript/data/vlm_trial_ledger.json; scripts/audit_vlm_reconciliation.py; manuscript/TRIX_SUPPLEMENT.tex}

\response{R4.10}{BATTERY threshold and thermal interpretation}{Reconcile the 20-N force threshold, thermal criterion and inconsistent table/text descriptions.}
The revised BATTERY model distinguishes a 15-N preventive deformation onset,
the 20-N flagged vertical-force boundary, a 15-N lateral boundary, accumulated
deformation/internal-short state and the diagnostic temperature threshold
$T>60^\circ$C. These are different quantities and are now listed separately
with their outcome roles. A current force limit cannot remove thermal energy
or an internal short accumulated earlier; the delayed-state limitation remains.
Exact historical temperature/time examples without supporting traces have
been removed.

The matched-specification comparison supplies both the box and TRiX arms
with the same preventive 15-N target and margin construction. Both attain
100\% safe completion, so the earlier difference cannot be attributed to
projection form. Separately, the aggregate audit corrected the Stage~1
SAC/TRiX value in \msitem{tab:submitted-4} to 20.0\% after removing pooled preventive-arm shards.
The exact-arm audit in \path{docs/battery_aggregation_correction.md} and
\path{verification/policy_record_linkage.json} in the evidence archive identifies
the \texttt{stage1/BATTERY/sac/trix} group and its shards. The physical
specification correction and bookkeeping correction are documented separately,
and the original records remain preserved.
\locations{\mssec{2.3}; \mssec{3}; \mssec{4.3}; \mssec{S2.3}; \mssec{S7.1}; \mssec{S9.4}; \mssec{S13.1}; \msitem{tab:submitted-4}; \msitem{tab:submitted-8}}
\evidence{envs/battery_env_v2.py; docs/battery_aggregation_correction.md; archive/verification/policy_record_linkage.json; manuscript/data/seed_statistics_summary.json; manuscript/TRIX_SUPPLEMENT.tex}

\response{R4.11}{Simulation and control rates}{Resolve the inconsistent 240-Hz and 100-Hz descriptions.}
All six revised analytical environments import $\Delta t=1/240$ s from
\texttt{benchmark/disasm\_bench.py}; each action step advances approximately
4.167 ms of simulated time. SCREW uses 40 internal integration substeps;
PCB, SNAP, CRANK and BATTERY use ten, while PRY updates at the outer step.
Supplementary Section~S4.2 now states these values explicitly and withdraws the submitted
100-Hz simulation statement. The approximately 200-Hz physical B601-RS
controller is described separately. Simulated-time step size is not a
measurement of wall-clock governor throughput or a real-time deadline guarantee.
\locations{\mssec{4.1}; \mssec{4.7}; \mssec{4.8}; \mssec{S4.2}; \mssec{S8}; \msitem{para:simulation-timestep}}
\evidence{benchmark/disasm_bench.py; envs/screw_env_v3.py; envs/pcb_env_v2.py; envs/snap_env_v2.py; envs/crank_env_v2.py; envs/battery_env_v2.py; envs/pry_env_v2.py; manuscript/TRIX_SUPPLEMENT.tex}

\response{R4.12}{Figure~4 and Table~4 use different experiments}{Resolve the mismatch between the submitted learning curves and headline results.}
The old Figure~4 has been replaced. The revised figure shows PCB's
post-hoc/filter-aware comparison and the Stage~2 SNAP, CRANK and PRY
comparisons, using the frozen production records reported in the revision.
Its caption explicitly distinguishes the regimes. The corresponding
post-hoc and filter-aware summaries are now \msitem{tab:submitted-4}
and \msitem{tab:stage2-key}, respectively. Main Table~4 gives the key contrasts.
Figure~4 was regenerated from the retained seed-level data with the supplied
replacement plotting script. All 80 plotted values are checked against their
recorded source cells before plotting. No curve from a simplified environment is used to explain another
environment's absolute rates.
\locations{\mssec{2.3}; \mssec{S9}; \msitem{tab:submitted-4}; \msitem{tab:stage2-key}; \msitem{fig:seed-distributions}}
\evidence{manuscript/figures/figure4_seed_data.csv; scripts/build_revision_figure4.py; manuscript/figures/figure4_rebuild.json; results/stage1/stage1_records.json; results/stage2/stage2_records.json; manuscript/TRIX_SUPPLEMENT.tex}

\response{R4.13}{PCB residual violations and representation analysis}{Explain the reported PCB failures instead of attributing them generically to curvature or linearity.}
We replaced the submitted interpretation with a decomposition of the
implemented constraint representations and retained the specification caveat.
In the archived PCB proposal study, clipping yields 2.9\% feasible commands;
the learned action-affine model yields 0.0\%; nonlinear one-step correction
yields 78.0\% and iteration yields 92.3\%; analytical oracle and exact
projection each yield 100\%. This is one fitting run with 20,000 training
samples and 2,000 held-out proposals, not a trained-policy safety ranking.
Increasing fitting samples from 500 to 20,000 does not resolve the affine
model's zero-feasibility result in that experiment.

The production PCB comparison also changes two things: the historical box
uses a nominal 0.100-N\,m per-axis bound, while TRiX uses a robust coupled
radius of about 0.087087 N\,m, a 12.913\% reduction. The original gap
therefore does not isolate box-versus-disk geometry. Holding the box filter
fixed across regimes, SAC safe completion changes from 0.0\% post hoc to
100.0\% with filter-aware training. The general representation analysis
uses the separate procedural benchmark with its disclosed fitting and
partition design. This identifies the roles of learned approximation,
constraint specification and policy adaptation without preserving the
unsupported curvature explanation of the old 44\% figure.
\locations{\mssec{2.3}; \mssec{2.4}; \mssec{4.3}; \mssec{4.4}; \mssec{S9.3}; \mssec{S10.2}; \mssec{S10.3}; \msitem{tab:procedural-geometry}}
\evidence{results/safelayer_decomposition.json; baselines/pcb_filters.py; docs/provenance.md; results/geometry_bench_heldout.json; manuscript/TRIX_SUPPLEMENT.tex}

\end{document}
